\documentclass[11pt]{article}
\usepackage[margin=1in]{geometry}
\usepackage{graphicx,booktabs,amsmath,float,hyperref,caption}
\graphicspath{{../results/plots/}}
\title{Reinforcement Learning for Dynamic Food Delivery Dispatch\\\large Implementation and Experimental Report}
\author{Anshit Agarwal \and Namya Dhingra}
\date{October 2026}

\begin{document}
\maketitle

\section{Problem and approach}
A central dispatcher must decide, at every moment, which delivery partner serves which order and
whether to wait. We formulate this as an MDP $(S,A,P,R,\gamma)$ over a 10$\times$10 city grid and
compare three heuristics (nearest driver, most urgent order, minimum estimated delivery time) with
two learned policies (DQN, PPO) \emph{in the same simulator, on the same seeds}. Every episode lasts
exactly $T=3$ simulated minutes (180 one-second ticks), independent of the configuration or the
algorithm.

\section{Environment and MDP}
\paragraph{Simulator.} Restaurants, customers, drivers and orders live on a continuous grid with Manhattan
travel (L-shaped routes). Orders arrive by a Poisson process (plus a small backlog at $t=0$), have a hidden
true and an observable estimated preparation time, and a deadline. Order lifecycle:
Created $\to$ Preparing $\to$ Ready $\to$ Assigned $\to$ Picked up $\to$ Delivered; a driver can be
assigned while food is still being prepared and then waits at the restaurant. All randomness is
pre-sampled from the seed into a scenario trace so that every strategy faces an identical world.

\paragraph{State.} A fixed 481-dimensional vector: 12 driver slots (position, availability, load, status,
time until free), 16 order slots (restaurant/customer position, age, estimated prep remaining, ready flag,
time to deadline, restaurant load), 12$\times$16 estimated completion times for each valid driver--order pair,
and 9 global features. Hidden quantities are never exposed.

\paragraph{Actions.} $\text{assign}(d_i,o_j)$ for $12\times16$ slots plus $\text{wait}$ (193 actions), with an
action mask; invalid actions are treated as waiting and penalised.

\paragraph{Reward.}
\[
R_t = w_1 D_t - w_2 W_t - w_3 L_t - w_4 T_t + w_5 B_t - w_6 I_t,
\]
with deliveries $D_t$, order-minutes waiting $W_t$, order-minutes overdue $L_t$, distance $T_t$, batching
savings $B_t$ and idle-driver-minutes $I_t$, plus a penalty for each order undelivered at $t=180$. The weights
(10, 2, 4, 0.1, 0.5, 1; undelivered 5) were fixed before the final evaluation.

\section{Methods}
\textbf{Baselines.} B1 serves the oldest order with the nearest driver; B2 serves the order with the least
slack to its deadline with the nearest feasible driver; B3 greedily picks the feasible (driver, order) pair
with the smallest estimated delivery time.

\textbf{DQN.} Double DQN with action masking, uniform replay, a hard-updated target network,
$\epsilon$-greedy exploration among valid actions and Huber loss.

\textbf{PPO.} Separate policy and value networks, masked categorical policy, GAE($\lambda{=}0.95$),
clipped surrogate objective, entropy bonus, parallel rollouts.

Both use a pairwise network that scores each (driver, order) pair with shared weights and a separate head for the wait action.

\section{Experimental protocol}
Five scenarios (Low/Medium/High demand, Driver shortage/surplus). Training samples random scenarios; model
selection uses 20 validation episodes; final results use 100 held-out test seeds per scenario (disjoint
from training and validation seeds). RL results are averaged over three independent training seeds.
Paired Wilcoxon tests compare each RL agent with each baseline on identical seeds.

\section{Results}
% Placeholder: fill in after running the final evaluation (see README "Quick start").
% Tables:  \input{../results/tables/summary_final_medium_demand.tex}  (one file per scenario)
% Figures: \includegraphics[width=.48\textwidth]{01_avg_delivery_time.pdf} ... (results/plots/)
Results are generated by \texttt{python -m evaluation.comparison} and
\texttt{python -m evaluation.visualization}; paste the resulting tables and discuss the trade-offs here.


\section{Limitations}
The simulator is deliberately small: time is compressed so that a delivery takes about a minute, there is
one pickup per trip unless batching is enabled, travel follows Manhattan paths without road networks, and
the RL agents were trained without batching.

\end{document}
